\sscp{Reduction of medial *CC clusters}{13-02-01}
\citet[246]{Blust1993} saw “the reduction of consonant clusters
in the reflexes of reduplicated monosyllables” as a phonological
innovation that distinguishes “PMP from PCEMP”.
He further explained, 

\begin{quote}
		“So far as is known,
		all CEMP languages have simplified medial clusters in the reflexes
		of PMP reduplicated monosyllables unless the cluster consisted
		of a nasal followed by a stop or fricative,
		in which case the nasal assimilated to the place of articulation
		of the obstruent, but was not lost: for example,
		PMP *bukbuk > PCEMP *bubuk ‘wood weevil’,
		PMP *ñamñam > PCEMP *ñañam ‘tasty,
		delicious’, PMP *mekmek > PCEMP *memek ‘crumbs’,
		but PMP *demdem > PCEMP *dendem ‘dark’.
		Some WMP languages (as Malay) have simplified PMP reduplicative
		clusters in much the same way,
		but the universality of this change in the CEMP languages is
		most simply explained as the product of a single innovation in
		a language ancestral to the entire group.” \citep[246]{Blust1993}
\end{quote}

\citet[134]{DonohueGrimes2008} pointed out the pattern is both
widespread and expected, “The process of cluster reduction is widely attested
in the CMP area,
but is also an expected feature of any language that approximates
a CV syllable structure.
Outside the CEMP area,
in Southeast Sulawesi, for instance,
we find that roots of the original form
*C\sub{1}V\sub{1}C\sub{2}C\sub{3}V\sub{2}C\sub{4} are realised
as C\sub{1}V\sub{1}C\sub{3}V\sub{2},
showing an even more complete adherence to the cross-linguistically
unmarked CV template, in examples such as
*tuktuk > Tukang Besi \it{tutu} ‘pound’ (Donohue 1999) [\citealp{Donohue1999a}].”
A number of languages in the southern Philippines,
western Indonesia and Sulawesi show similar reduction of the
medial **CC of PMP *tuktuk ‘pound’ (e.g.
Maranao \it{totok}, Iban \it{tutok},
Malay \it{tutok}, Kulisusu \it{mon-cucu},
Wolio \it{tutuk-i}; data from \citealp{BlustTrussel2020}).
This is only one of many examples of how widespread the reduction occurs.
Such reduction of medial *CC clusters is not exclusive to the CEMP region.

\citet[222]{Blust2013} also notes patterns where the coda of
the first syllable is lost,
but the coda of the final syllable is retained.
“PAN permitted heterorganic consonant clusters in reduplicated
monosyllables such as *buCbuC ‘pluck,
pull out’, *kiskis ‘scrape’,
or *tuktuk ‘knock, pound, beat’.
These have been retained in many languages,
and in several of the Formosan languages they are the only clusters allowed.
All Central-Eastern Malayo-Polynesian languages
and some languages in western Indonesia have reduced such clusters,
with the exception that nasal-obstruent sequences are sometimes
retained: *seksek > Manggarai \it{cəcək} ‘stuff,
cram in’, *bejbej > Kambera \it{búburu} ‘bundle’,
*bukbuk > Tetun \it{fuhuk} ‘wood weevil’,
*sepsep > Javanese \it{səsəp} ‘suck’,
*buŋbuŋ > \it{wuwuŋ} ‘ridge of the roof’,
*bukbuk > Malay \it{bubok} ‘wood weevil’,
but *diŋdiŋ > \it{dindiŋ} ‘wall’.
As these examples illustrate,
the loss of syllable-final consonants in medial position is independent
of the loss of word-final consonants.” Again,
this is to be expected in languages whose phonotactics require
words to be shaped CVCVC rather than CVCCVC.

There are several logical possibilities of what can happen to
medial CC:

\begin{itemize}
		\item	The first consonant is preserved; the second is lost: C\sub{1}C\sub{2} → C\sub{1}
		\item	The second consonant is preserved; the first is lost: C\sub{1}C\sub{2} → C\sub{2}
		\item	The two consonants merge to become a different sound: C\sub{1}C\sub{2} → C\sub{3}
		\item	Both consonants are preserved
					(either unchanged or with changes to one or both of the consonants:
					C\sub{1}C\sub{2} → C\sub{1}C\sub{2} or C\sub{1}′C\sub{2}′)
\end{itemize}

All four patterns are found in the target region,
particularly if C\sub{1} is a nasal *NC.
We note that *NC clusters are not included in Blust’s claims,
but we include those in our broader discussion of all medial
clusters as background for the more specific claims that are
limited to medial *CC clusters in reduplicated monosyllables (excluding *NC clusters).
Note Blust’s claim accounts only for the second pattern above,
whereas the variety summarised above found in the region means
that the full historical *CC clusters must (often) be reconstructed
to account for all the data in the region.

\citet[54]{Blust2008} provides the following examples of what
happened to historical medial *C\sub{1}C\sub{2} → C\sub{2} in
Kambera, a \tsc{Bima-Lembata} language spoken on Sumba.
We supplement his Kambera examples in \trf{13-05} with data from
other languages of Sumba.
These show the reduction of medial *CC clusters not involving nasals.
These examples focus on the broader issue of medial clusters,
and are not restricted to reduplicated monosyllables.
%(Note: Kodi orthographic ‹gh› = velar fricative [ɣ]; Anakalang
%and Mamboru orthographic ‹ḅ› are implosive [ɓ].

\begin{table}
\tcp{Medial *C\sub{1}C\sub{2} > C\sub{2} in languages of Sumba}{13-05}
\begin{adjustbox}{width=\textwidth}
	\begin{threeparttable}\st{2pt}
	\begin{tabular}{lllllll}\lsptoprule
local gl.	&	  porridge	&	  bundle	&	rudder, row	&	 penetrate	&	puff, blow on	&	beat, knock, pound\\
PMP	&	*buR\tbr{b}uR	&	*bəj\tbr{b}əj	&	*bəR\tbr{s}ay	&	*suk\tbr{s}uk	&	*put\tbr{p}ut	&	*tuk\tbr{t}uk\\	\midrule
Kambera	&	{\hr}bu\tbr{b}ur-u\su{†}	&	{\hr}bu\tbr{b}ur-u	&	‹bú\tbr{h}i›	&	{\hr}hu\tbr{h}uk-u	&	{\hr}pu\tbr{p}ut-u	&	{\hr}tu\tbr{t}uk-u\\
Mangili	&	{\hr}bu\tbr{b}ur-u	&		&	‹bú\tbr{h}i›	&	{\hr}hu\tbr{h}uk-u	&	{\hr}pu\tbr{p}ut-u	&	{\hr}tu\tbr{t}uk-u\\
Lewa	&	{\hr}bu\tbr{b}ur-u	&		&	{\hr}bo\tbr{h}u	&	{\hr}hu\tbr{h}uk-u	&	{\hr}pu\tbr{p}ut-u	&	{\hr}to\tbr{t}uk-u\\
Anakalang	&	{\hr}bu\tbr{b}ur-u	&		&	{\hr}ɓo\tbr{s}i	&	{\hr}su\tbr{s}uk-a,&	{\hr}pu\tbr{p}ut-u	&	{\hr}to\tbr{t}uk-u\\
					&											&		&									&{\hr}si\tbr{s}ik-u	&		&	\\
Mamboru	&	{\hr}ba\tbr{b}ur-a	&		&	{\hr}ɓo\tbr{s}i	&	{\hr}su\tbr{s}uk-a	&	{\hr}pu\tbr{p}ut-a	&	‹tú\tbr{t}u›\\
Weyewa	&	{\hr}bu\tbr{b}ur-a	&	‹wò\tbr{b}ola›	&	{\hr}bo\tbr{z}e	&		&	{\hr}po\tbr{w}i	&	{\hr}katu\tbr{t}uk-a\\
Kodi	&	{\hr}bu\tbr{b}ur-o	&	{\hr}ɣo\tbr{b}ːol-o	&	{\hr}bo\tbr{h}e	&	{\hr}hu\tbr{h}uk-o	&	{\hr}pu\tbr{p}ut-o	&	{\hr}tu\tbr{t}uk-o\\
		\lspbottomrule
	\end{tabular}
		\begin{tablenotes}\tnsize
			\item[†] {Most final vowels on trisyllables are extra-metrical paragogic (epenthetic) vowels.}
		\end{tablenotes}
	\end{threeparttable}
\end{adjustbox}
\end{table}
